% How Many Ways to Walk Forward?  Main manuscript.
% Build: latexmk -pdf main.tex   (numbers.tex is generated by make_numbers.py)
\documentclass[11pt,a4paper]{article}

% ------------------------------------------------------------------ layout
% A4 page with a 6.5in text width, so that the figures, rendered at 624px = 6.5in, are included at
% 100% of their design size. Tables are typeset by LaTeX from ../tables/tex.
\usepackage[a4paper,textwidth=6.5in,top=1in,bottom=1.1in]{geometry}
\usepackage[T1]{fontenc}
\usepackage{stix2}                            % text and math
\usepackage[scaled=0.95]{sourcesanspro}
\usepackage[final]{microtype}
\usepackage{setspace}\setstretch{1.12}        % use \doublespacing for submission if required
\usepackage{amsmath,amsthm}
\usepackage{graphicx}
\usepackage{booktabs,tabularx}
\usepackage[section]{placeins}
\usepackage{enumitem}
\usepackage[labelfont={bf,sf},font=small,labelsep=period,skip=6pt]{caption}
\captionsetup[table]{position=top,justification=raggedright,singlelinecheck=false}
\usepackage{etoolbox}
\AtBeginEnvironment{table}{\small\setstretch{1.0}\setlength{\tabcolsep}{4pt}}
\usepackage[dvipsnames]{xcolor}
\usepackage{titlesec}
\usepackage{authblk}
\usepackage[ruled,linesnumbered,vlined]{algorithm2e}
\usepackage[style=authoryear,backend=biber,maxcitenames=2,mincitenames=1,maxbibnames=7,minbibnames=6,
            uniquelist=false,uniquename=false,giveninits=true,doi=true,url=false,eprint=true,
            isbn=false,dashed=false,sortcites=true]{biblatex}
\addbibresource{references.bib}
\DeclareFieldFormat[article]{title}{#1}
\renewbibmacro{in:}{\ifentrytype{article}{}{\printtext{\bibstring{in}\intitlepunct}}}
\AtEveryBibitem{\ifentrytype{misc}{}{\clearfield{url}}}
\usepackage[colorlinks=true,linkcolor=MidnightBlue,citecolor=MidnightBlue,urlcolor=MidnightBlue]{hyperref}
\hypersetup{pdftitle={How Many Ways to Walk Forward? Design Snooping in Trading-Rule Backtests},
            pdfauthor={Carlos A. Castro, Rodrigo A. Castro, Paula M. Almonacid}}
\usepackage{orcidlink}                       % after hyperref

\titleformat{\section}{\normalfont\large\bfseries}{\thesection.}{0.5em}{}
\titleformat{\subsection}{\normalfont\normalsize\bfseries}{\thesubsection.}{0.5em}{}
\titleformat{\paragraph}[runin]{\normalfont\normalsize\itshape}{}{}{}[.]
\titlespacing*{\paragraph}{0pt}{6pt}{0.6em}
\setlist{itemsep=2pt,topsep=4pt}
\renewcommand\Affilfont{\small}
\setlength{\affilsep}{0.6em}
% let pages hold more than one float and little text
\renewcommand{\topfraction}{0.9}\renewcommand{\bottomfraction}{0.6}
\renewcommand{\textfraction}{0.07}\renewcommand{\floatpagefraction}{0.75}
\setcounter{topnumber}{3}\setcounter{totalnumber}{4}

\newtheorem{definition}{Definition}

\newcommand{\E}{\mathbb{E}}
\newcommand{\figpath}{../figures/pdf}
\newcommand{\tabpath}{../tables/tex}
% Generated by make_numbers.py; do not edit.
\newcommand{\NumDate}{November 2026}
\newcommand{\NumRules}{351}
\newcommand{\NumRulesFilter}{93}
\newcommand{\NumRulesMA}{116}
\newcommand{\NumRulesSR}{70}
\newcommand{\NumRulesCB}{72}
\newcommand{\NumBurnIn}{250}
\newcommand{\NumGarchOmega}{\ensuremath{1.22\times10^{-6}}}
\newcommand{\NumGarchAlpha}{0.097}
\newcommand{\NumGarchBeta}{0.896}
\newcommand{\NumGarchNu}{6.86}
\newcommand{\NumNullReps}{5,000}
\newcommand{\NumAltReps}{2,000}
\newcommand{\NumT}{5,292}
\newcommand{\NumNoos}{4,013}
\newcommand{\NumPhiList}{0.02, 0.04, 0.06, 0.08, 0.10, 0.12, 0.15}
\newcommand{\NumMuList}{10\%, 20\%, 30\%, 40\%, 50\%}
\newcommand{\NumFullbootB}{299}
\newcommand{\NumFullbootReps}{1,200}
\newcommand{\NumFullbootMax}{3.1}
\newcommand{\NumFullbootAvg}{3.3}
\newcommand{\NumFullbootMaxIid}{4.7}
\newcommand{\NumFullbootAvgIid}{4.4}
\newcommand{\NumSubsets}{25}
\newcommand{\NumSizeSingleMin}{4.4}
\newcommand{\NumSizeSingleMax}{5.2}
\newcommand{\NumSizeSingle}{4.8}
\newcommand{\NumSizeBestTwo}{7.2}
\newcommand{\NumSizeBestEight}{12.5}
\newcommand{\NumSizeBestNaive}{20.5}
\newcommand{\NumSizeBestNaiveRound}{4.1}
\newcommand{\NumSizeBonf}{0.6}
\newcommand{\NumSizeDSR}{2.9}
\newcommand{\NumSizeBonfEff}{14.6}
\newcommand{\NumSizeDSREff}{17.3}
\newcommand{\NumSizeMaxT}{3.7}
\newcommand{\NumSizeAvgT}{5.0}
\newcommand{\NumSizeRulesNaive}{40.1}
\newcommand{\NumSizeRulesRC}{3.3}
\newcommand{\NumISann}{13.3}
\newcommand{\NumOOSann}{0.0}
\newcommand{\NumMeanCorrT}{0.79}
\newcommand{\NumMinCorrT}{0.71}
\newcommand{\NumNeffT}{1.6}
\newcommand{\NumNindep}{4.5}
\newcommand{\NumNindepRound}{4}
\newcommand{\NumMeanCorrD}{0.79}
\newcommand{\NumNeffD}{1.6}
\newcommand{\NumMCdistinct}{19}
\newcommand{\NumMCtiedShare}{6}
\newcommand{\NumPowerPhi}{0.10}
\newcommand{\NumPowerAvg}{39.1}
\newcommand{\NumPowerMax}{41.2}
\newcommand{\NumPowerMedian}{33.1}
\newcommand{\NumPowerSingleMin}{16}
\newcommand{\NumPowerSingleMax}{55}
\newcommand{\NumPowerBest}{54.6}
\newcommand{\NumPowerExpanding}{54.2}
\newcommand{\NumPowerNaiveAdj}{42.1}
\newcommand{\NumPowerNaive}{72.0}
\newcommand{\NumPowerBonfEff}{63.4}
\newcommand{\NumPowerDSREff}{69.0}
\newcommand{\NumPowerMu}{30\%}
\newcommand{\NumPowerAvgMu}{40.7}
\newcommand{\NumPowerMaxMu}{39.9}
\newcommand{\NumPowerMedianMu}{35.8}
\newcommand{\NumPowerExpandingMu}{39.9}
\newcommand{\NumPowerNaiveAdjMu}{41.4}
\newcommand{\NumPowerGapPoints}{15}
\newcommand{\NumPowerGapBest}{15}
\newcommand{\NumSizeBonfIndep}{6.3}
\newcommand{\NumFullbootSE}{0.6}
\newcommand{\NumFullbootRepsIid}{1,200}
\newcommand{\NumSizeSingleCash}{5.0}
\newcommand{\NumSizeBestNaiveCash}{47.8}
\newcommand{\NumSizeMaxTCash}{3.2}
\newcommand{\NumSizeAvgTCash}{6.0}
\newcommand{\NumSizeBestNaiveCashIid}{46.9}
\newcommand{\NumSizeMaxTCashIid}{2.9}
\newcommand{\NumSizeAvgTCashIid}{4.8}
\newcommand{\NumSizeRobustMin}{2.9}
\newcommand{\NumSizeRobustMax}{6.0}
\newcommand{\NumMeanCorrTCash}{0.38}
\newcommand{\NumMinCorrTCash}{0.17}
\newcommand{\NumMeanCorrCash}{0.37}
\newcommand{\NumNindepCash}{12.7}
\newcommand{\NumNindepCashRound}{13}
\newcommand{\NumNeffLiji}{12.1}
\newcommand{\NumNeffGalwey}{11.9}
\newcommand{\NumNeffNyholt}{21.6}
\newcommand{\NumNeffNyholtRound}{22}
\newcommand{\NumSizeBonfLiji}{2.8}
\newcommand{\NumSizeBonfGalwey}{2.9}
\newcommand{\NumSizeBonfNyholt}{1.4}
\newcommand{\NumPremium}{6\%}
\newcommand{\NumPremiumReps}{2,000}
\newcommand{\NumSizeSinglePremium}{0.3}
\newcommand{\NumSizeBestNaivePremium}{3.1}
\newcommand{\NumCapPct}{50\%}
\newcommand{\NumCappedDays}{69}
\newcommand{\NumCappedReps}{38}
\newcommand{\NumFullbootMaxCash}{4.0}
\newcommand{\NumFullbootAvgCash}{4.5}
\newcommand{\NumFullbootMaxCashIid}{4.2}
\newcommand{\NumFullbootAvgCashIid}{3.8}
\newcommand{\NumFullbootPowerReps}{1,000}
\newcommand{\NumFullbootPowerMax}{42.8}
\newcommand{\NumFullbootPowerAvg}{39.0}
\newcommand{\NumFullbootPowerMaxMu}{38.0}
\newcommand{\NumFullbootPowerAvgMu}{38.6}
\newcommand{\NumUSoosFirst}{1930}
\newcommand{\NumUSgrossSig}{36}
\newcommand{\NumUSgrossBestT}{3.65}
\newcommand{\NumUSgrossMinT}{$-$0.04}
\newcommand{\NumUSgrossPmax}{0.002}
\newcommand{\NumUSgrossPavg}{0.017}
\newcommand{\NumUSnetSig}{0}
\newcommand{\NumUSnetBestT}{$-$0.13}
\newcommand{\NumUSnetMeanT}{$-$2.38}
\newcommand{\NumUSnetPmax}{0.843}
\newcommand{\NumUSnetPavg}{0.996}
\newcommand{\NumUSmeanCorr}{0.84}
\newcommand{\NumUSneff}{1.4}
\newcommand{\NumUSdistinctLow}{19}
\newcommand{\NumUStopRule}{F(0.005,\allowbreak e=5)}
\newcommand{\NumUStopShare}{26}
\newcommand{\NumUStopThreeShare}{58}
\newcommand{\NumUStiedWindows}{906}
\newcommand{\NumUSwindows}{26,474}
\newcommand{\NumUStiedShare}{3.4}
\newcommand{\NumUSturnover}{53}
\newcommand{\NumUSturnoverUnits}{107}
\newcommand{\NumUSmaxTiedDesign}{R252/h21/g21}
\newcommand{\NumUSmaxTied}{116}
\newcommand{\NumUSmaxTiedBlocks}{1194}
\newcommand{\NumUSmaxTiedDesc}{rolling 252-day window, 21-day test blocks, 21-day gap}
\newcommand{\NumNoosShareUS}{16}
\newcommand{\NumNoosYears}{16}
\newcommand{\NumACone}{0.08}
\newcommand{\NumACtwo}{0.20}
\newcommand{\NumACthree}{$-$0.06}
\newcommand{\NumACpre}{0.11}
\newcommand{\NumACpost}{$-$0.03}
\newcommand{\NumAPac}{0.12}
\newcommand{\NumAPnetPavg}{0.032}
\newcommand{\NumAPnetPmax}{0.015}
\newcommand{\NumAPgrossSig}{55}
\newcommand{\NumDVgrossSig}{47}
\newcommand{\NumAPnetPBO}{27}
\newcommand{\NumAPnetSig}{37}
\newcommand{\NumAPnetStepm}{12}
\newcommand{\NumUSgrossSharePreFutures}{87}
\newcommand{\NumUSnetRelPreFutures}{1.6}
\newcommand{\NumUSnetRelEndEighties}{1.6}
\newcommand{\NumUSnetRelEnd}{$-$5.0}
\newcommand{\NumUSnetPBO}{12}
\newcommand{\NumRegionsZero}{3}
\newcommand{\NumEraOneSig}{2}
\newcommand{\NumEraOneNetSig}{0}
\newcommand{\NumEraTwoNetStepm}{38}
\newcommand{\NumEraOnePmax}{0.158}
\newcommand{\NumEraOnePavg}{0.216}
\newcommand{\NumEraTwoSig}{56}
\newcommand{\NumEraTwoNetSig}{51}
\newcommand{\NumEraTwoNetPavg}{0.007}
\newcommand{\NumEraThreeSig}{0}
\newcommand{\NumEraThreeMaxT}{$-$0.71}
\newcommand{\NumUSlagSig}{0}
\newcommand{\NumUSlagBestT}{$-$1.83}
\newcommand{\NumAPlagSig}{0}
\newcommand{\NumAPlagBestT}{0.45}
\newcommand{\NumAPgrossBestT}{5.64}
\newcommand{\NumAPlagNetSig}{0}
\newcommand{\NumDVlagSig}{0}
\newcommand{\NumLagMaxSig}{0}
\newcommand{\NumLagBestT}{0.45}
\newcommand{\NumLagMinPmax}{0.702}
\newcommand{\NumLagSeries}{7}
\newcommand{\NumUSlagTurnover}{10}
\newcommand{\NumEraTwoLagSig}{0}
\newcommand{\NumEraTwoLagMaxT}{0.41}
\newcommand{\NumEraTwoBreakeven}{19}
\newcommand{\NumEraTwoMaxT}{6.49}
\newcommand{\NumCostBp}{10}
\newcommand{\NumOldRules}{211}
\newcommand{\NumOldSizeBestNaive}{16.2}
\newcommand{\NumUSclassicRule}{F(0.005)}
\newcommand{\NumUSclassicT}{2.82}
\newcommand{\NumUSclassicPrc}{0.059}


% ------------------------------------------------------------------ front matter
\title{\textbf{How Many Ways to Walk Forward?}\\[4pt]
\large Design Snooping in Trading-Rule Backtests}
\author[1]{Carlos A. Castro\,\orcidlink{0000-0001-8748-9918}}
\author[2,*]{Rodrigo A. Castro\,\orcidlink{0000-0003-0010-7775}}
\author[1]{Paula M. Almonacid\,\orcidlink{0000-0001-6008-1728}}
\affil[1]{Universidad EAFIT, School of Economics and Finance, Medell\'in, Colombia.}
\affil[2]{Universidad de Valpara\'iso, Institute of Mathematics, Valpara\'iso, Chile.}
\affil[*]{Corresponding author: \href{mailto:rodrigo.castrom@uv.cl}{rodrigo.castrom@uv.cl}}
\date{\vspace{-1.2em}}

\begin{document}
\maketitle
\thispagestyle{empty}

\begin{abstract}
\noindent Walk-forward analysis selects a trading rule on past data and trades it on the data that
follow, so the rule is never judged on the data that chose it. That removes the classic form of data
snooping but not a second one: the analysis has settings of its own (training window, test block,
gap), and a researcher who reports the setting that worked best has snooped again. We measure this
design snooping. In \NumNullReps{} simulated markets with GARCH
volatility and no predictability, 56 walk-forward designs are applied to \NumRules{} trading rules. A
design fixed in advance rejects at close to the nominal 5\%. The best of the 56 rejects in
\NumSizeBestNaive\% of markets when the benchmark is buy-and-hold and in \NumSizeBestNaiveCash\% when
it is cash, because subtracting a common benchmark return makes the designs' test statistics more alike
(mean pairwise correlation \NumMeanCorrT{} against \NumMeanCorrTCash). Two bootstrap tests across
designs, a Reality Check on the maximum and a test of the average $t$-statistic, keep the false-positive
rate between \NumSizeRobustMin\% and \NumSizeRobustMax\% under both benchmarks without a count of
trials, and against buy-and-hold they have the power of the median design. On US market returns from
\NumUSoosFirst{} to 2026, no design beats buy-and-hold after costs of \NumCostBp{} basis points per unit
traded. Before costs \NumUSgrossSig{} of 56 do, but none does when trades are executed one day after the
signal, which points to nonsynchronous index prices, not tradable predictability. Walk-forward results
should be reported for the full grid of designs, with the benchmark and a design-robust test.

\medskip\noindent\textbf{Keywords:} walk-forward analysis; backtest overfitting; data snooping;
out-of-sample evaluation; multiple testing; technical trading rules; bootstrap.

\noindent\textbf{JEL classification:} C12; C15; C22; C53; G11; G14.
\end{abstract}

\newpage
\setcounter{page}{1}

% =====================================================================================
\section{Introduction}\label{sec:intro}

Because it is easy to try thousands of trading rules on the same few decades of prices, a backtest that
looks good is often good only by chance \parencite{lo1990,white2000,bailey2014notices,harvey2016}.
Walk-forward analysis is the practitioner's standard answer \parencite{pardo2008}: the analyst chooses
the rule on a training window, trades it on the next, unseen test block, moves both windows forward and
repeats. The record then contains only returns the strategy could have earned in real time, and the same
idea underlies rolling-origin evaluation in forecasting \parencite{tashman2000,hyndman2021} and
prequential evaluation in machine learning \parencite{gama2013}.

Walk-forward analysis has settings of its own: whether the training window keeps a fixed length, grows or
uses all past data; how long it is; how long each test block is; and whether a gap separates the two.
There is no consensus on these choices and results change with them. \textcite{rossi2012} showed in
forecast evaluation that the window size can be mined, and \textcite{hansen2012split} showed the same
for the sample split. A researcher who tries several walk-forward designs and reports the one that worked
has snooped at the level of the evaluation design rather than the rule; practitioner guidance treats the
settings as parameters to be optimised \parencite{pardo2008}, and comparisons of validation schemes treat
walk-forward analysis as one fixed procedure \parencite{arian2024}.

This paper measures that design snooping and shows how to remove it. The yardstick is the false-positive
rate: the share of markets with no predictability in which a procedure declares that the strategy beats a
passive benchmark.

\begin{enumerate}[leftmargin=*]
\item \textbf{The size of design snooping, and what it depends on} (Sections~\ref{sec:family},
\ref{sec:design} and~\ref{sec:mc}). Rolling, growing and expanding schemes, with or without a gap, are
written as one family in which every design is scored on the same out-of-sample days and each day is
scored once. In \NumNullReps{} simulated markets with no predictability, a single design fixed in advance
has a false-positive rate between \NumSizeSingleMin\% and \NumSizeSingleMax\%. The best of $K$ designs
raises it with $K$: at $K=56$ to \NumSizeBestNaive\% when the benchmark is buy-and-hold and to
\NumSizeBestNaiveCash\% when it is cash. The difference has a simple source. The designs are dependent
for two reasons: they select among the same rules, and each is measured against the same benchmark
return, which enters every design's statistic. Against buy-and-hold the design statistics have a mean
pairwise correlation of \NumMeanCorrT{} and the best of 56 rejects as often as the best of about
\NumNindep{} independent tests; against cash the correlation is \NumMeanCorrTCash{} and the number is
about \NumNindepCashRound. The distortion is material under either benchmark for a grid that varies only
the windows, and it grows with whatever else the researcher varies.

\item \textbf{Design-robust tests} (Sections~\ref{sec:procedures} and~\ref{sec:mc}). Treating the designs
as a family of hypotheses and applying a bootstrap Reality Check \parencite{white2000} across them, or
testing the average $t$-statistic across designs as in \textcite{rossi2012}, keeps the false-positive
rate close to 5\%: between \NumSizeRobustMin\% and \NumSizeRobustMax\% across the two tests, the two
benchmarks and the two null processes. Both tests learn the dependence between designs from the data,
so they need no count of trials. Corrections that need one are only as good as the count: fed with the
nominal 56 or with four effective numbers taken from the literature, Bonferroni rejects in between
\NumSizeBonf\% and
\NumSizeBonfEff\% of null markets. Under planted predictability the two tests have the power of the
median single design and give up about \NumPowerGapBest{} percentage points against the best single
design, which no researcher could have known in advance.

\item \textbf{An illustration} (Section~\ref{sec:empirical}). On US market returns from \NumUSoosFirst{}
to 2026 and five regional indices since 1990, no design beats buy-and-hold after a cost of \NumCostBp{}
basis points per unit traded, with one regional exception. Before costs \NumUSgrossSig{} of 56 US designs
do and the design-robust tests reject, but not one design in any series remains significant when every
trade is executed one day after its signal. The profits come from first-order autocorrelation of the
index, which nonsynchronous prices create and no investor could trade at the closing price
\parencite{bessembinder1998,day2002,ready2002}. The illustration shows the tests at work; the empirical
findings themselves confirm earlier studies \parencite{sullivan1999,zakamulin2014,taylor2014}.
\end{enumerate}

% =====================================================================================
\section{Related work}\label{sec:related}

\paragraph{Out-of-sample evaluation and data snooping} Evaluating a strategy on data not used to build
it is the basic remedy for overfitting; with time series this leads to single splits, rolling or expanding
re-estimation and blocked cross-validation with buffers
\parencite{tashman2000,burman1994,racine2000,bergmeir2012,bergmeir2018,cerqueira2020,hewamalage2023,lopezdeprado2018},
and the choice between rolling and expanding schemes matters under instability
\parencite{west1996,clark2001,giacomini2006,pesaran2007,inoue2017,rossi2021}. When many rules are tested
on one history the significance of the best must account for the search: the Reality Check
\parencite{white2000}, its more powerful and step-down versions
\parencite{hansen2005,romano2005,hsu2010},
false-discovery-rate methods \parencite{benjamini1995,bajgrowicz2012,harvey2020raps}, the Deflated Sharpe
Ratio of the best of $N$ independent trials \parencite{bailey2014dsr}, with $N$ estimated by
clustering correlated trials \parencite{lopezdeprado2019}, the probability of backtest overfitting
\parencite{bailey2017} and stricter thresholds throughout empirical finance
\parencite{harvey2016,harvey2020jf}. Technical trading rules are the classic test case:
\textcite{brock1992} found that moving-average and range-breakout rules predicted Dow Jones returns;
\textcite{sullivan1999} showed that the best of nearly 8,000 rules survived the Reality Check up to 1986
but not afterwards; \textcite{bessembinder1998} showed that the forecast power of these rules weakens
when signals are executed a day late and does not cover trading costs, and \textcite{day2002} and
\textcite{ready2002} that much of the earlier profit reflects nonsynchronous prices and dividends rather
than tradable predictability; later studies find that
profits left major markets after the early 1990s \parencite{hsu2005,park2007,bajgrowicz2012,taylor2014,hsu2016};
and \textcite{zakamulin2014,zakamulin2017} gives the most thorough out-of-sample study of moving-average
timing on US data, with rolling and expanding windows, explicit split sensitivity, and profits
concentrated in early subperiods.

\paragraph{The evaluation design as a degree of freedom} The forecasting literature recognised early
that the sample split is a choice \parencite{inoue2005}. \textcite{hansen2012split} show that searching
over split points inflates the size of out-of-sample tests, and \textcite{rossi2012} propose tests that
remain valid when the window size is chosen after seeing the results, by taking the supremum or the
average of the statistic over window sizes. More broadly, undisclosed flexibility in how a study is
run inflates false positives \parencite{simmons2011}, and the ``non-standard errors'' and multiverse
literatures document that defensible design choices produce dispersion comparable to sampling error
\parencite{steegen2016,simonsohn2020,menkveld2024,soebhag2024,coqueret2023}. We bring these ideas to
walk-forward backtests, where the design has four or more settings and the strategy is re-selected in
every block; applied work that varies the windows searches over them rather than correcting for the
search \parencite[for example][]{mroziewicz2026}.

% =====================================================================================
\section{A unified walk-forward family}\label{sec:family}

\begin{figure}[tbp]
\centering
\includegraphics[width=\textwidth]{\figpath/fig1_design.pdf}
\caption{The three schemes of the walk-forward family and the optional gap. In every scheme the test
blocks are contiguous and tile the evaluation span, and every design starts testing on the same day.}
\label{fig:design}
\end{figure}

\subsection{Returns, positions and the quantity being tested}

Time is measured in trading days, $t=0,1,\dots,T-1$. We work with excess returns $e_t$, the market return
minus the one-month Treasury-bill rate.\footnote{Conventions: returns are daily simple returns; the
Sharpe ratio is the mean of daily excess returns divided by their standard deviation; tests are one-sided
at the 5\% level with $H_0$ that the strategy does not beat the benchmark; size is the rejection
probability under $H_0$ and power the rejection probability under a stated alternative; a basis point is
0.01\%; 252 trading days make a year.}
A trading rule $j$ takes a position $p_{j,t}\in\{-1,0,+1\}$ (short, flat, long) at the close of day $t$
using prices up to $t$, and earns $x_{j,t}=p_{j,t-1}\,e_t$ on day $t$ (the futures convention).

A strategy can be judged against two benchmarks. Against \emph{cash} the question is whether its mean
excess return is positive. Against \emph{buy-and-hold}, which earns $e_t$ every day, the question is
whether timing the market adds to holding it. Buy-and-hold is the benchmark of this paper, because a
timing rule on an equity index is of interest only if it beats holding the index. Cash appears in the
simulations because the size of design snooping turns out to depend on the benchmark.

\subsection{The family of designs}

A design $\theta$ consists of a test-block length $h$, an initial training length $L$, a per-block growth
$G\ge 0$, a gap $g\ge 0$ and a scheme flag for expanding windows. All designs share an origin $o$, the
first day any training window may use, and a first test day $s$.

\begin{definition}[Walk-forward family]\label{def:family}
For blocks $b=0,1,\dots,B-1$ with $B=\lceil (T-s)/h\rceil$, design $\theta$ uses the test block
\[
\mathcal{T}_b=\bigl[\,s+bh,\ \min\{s+(b+1)h,\,T\}\bigr)
\]
and the training window
\[
\mathcal{W}_b=\begin{cases}
\bigl[\max\{o,\ s+bh-g-L-Gb\},\ s+bh-g\bigr) & \text{rolling } (G=0) \text{ or growing-rolling } (G>0),\\[2pt]
\bigl[\,o,\ s+bh-g\bigr) & \text{expanding.}
\end{cases}
\]
\end{definition}

A rolling window keeps a fixed length $L$; a growing-rolling window starts at $L$ and gains $G$ days per
block; an expanding window uses everything since the origin (Figure~\ref{fig:design}). By construction
the test blocks of every design are disjoint and their union is $[s,T)$: every out-of-sample day is
scored exactly once, every training window ends at least $g+1$ days before its test block, and all
designs are scored on the same days, which is what makes them comparable and jointly testable. (Some
walk-forward implementations let test windows overlap; an average over overlapping windows is not the
return of any tradable strategy.) For one-day-ahead rules $g=0$ is already leak-free; a gap is needed
when the target spans several days \parencite{burman1994,racine2000,lopezdeprado2018}, and we include
it because practitioners use it, which makes it one more setting to vary.

\subsection{The walk-forward strategy}

In each block, design $\theta$ selects, among the $J$ rules, the one with the highest Sharpe ratio over
the training window and holds it, unchanged, through the test block (Algorithm~\ref{alg:wf}). The result
is a single daily series of out-of-sample returns $y^{\theta}_t$ for $t\in[s,T)$; nothing is tuned on the
test data and nothing is re-optimised within a block.

\begin{algorithm}[tbp]
\caption{Walk-forward meta-strategy for design $\theta$}\label{alg:wf}
\SetKwInOut{Input}{Input}\SetKwInOut{Output}{Output}
\DontPrintSemicolon
\Input{rule returns $x_{j,t}$ for $j=1,\dots,J$; design $\theta=(\text{scheme},L,G,h,g)$; origin $o$; first test day $s$}
\Output{out-of-sample returns $y^{\theta}_t$ for $t\in[s,T)$}
\For{$b=0,1,\dots,B-1$}{
  form $\mathcal{W}_b$ and $\mathcal{T}_b$ as in Definition~\ref{def:family}\;
  $\hat\jmath_b \leftarrow \arg\max_j \widehat{SR}_j(\mathcal{W}_b)$ \tcp*{in-sample Sharpe ratio; ties to the lowest $j$}
  $y^{\theta}_t \leftarrow x_{\hat\jmath_b,t}$ for all $t\in\mathcal{T}_b$ \tcp*{trade the chosen rule, untouched}
}
\end{algorithm}

% =====================================================================================
\section{Testing a strategy, and testing across designs}\label{sec:procedures}

\subsection{One design}\label{sec:onedesign}

For design $\theta$ the out-of-sample excess return over buy-and-hold is $d^{\theta}_t=y^{\theta}_t-e_t$,
or, net of costs, $d^{\theta}_t-c\,\lvert\Delta p^{\theta}_t\rvert$ with $c$ the cost per unit change in
position (a switch from long to short is two units). The null is $H_0^{\theta}:\E[d^{\theta}_t]\le 0$
and the statistic is
\[
t_{\theta}=\frac{\bar d^{\,\theta}}{\hat\sigma_{\theta}/\sqrt{n}},\qquad n=T-s,
\]
with $\hat\sigma^2_{\theta}$ the Newey--West long-run variance (Bartlett kernel, bandwidth
$\lfloor 4(n/100)^{2/9}\rfloor$; \cite{newey1987,newey1994}). The test rejects when $t_{\theta}>1.645$.
With cash as the benchmark, $d^{\theta}_t$ is replaced by $y^{\theta}_t$ throughout. Because each design
is scored once per day on a common span, this is an ordinary test on one time series. If excess returns
have zero mean and cannot be predicted from past prices, both $y^{\theta}_t$ and $d^{\theta}_t$ are
martingale differences whichever rule is selected, since positions use past prices only, and
Section~\ref{sec:mc} confirms that the size of the test is close to 5\%.

\subsection{Many designs}\label{sec:manydesigns}

Suppose a researcher runs $K$ designs with statistics $t_1,\dots,t_K$ and looks at all of them. Even if
no design has an edge, the largest of $K$ statistics is larger than a typical one. We compare three ways
of drawing a conclusion from the $K$ statistics, each at a nominal 5\% level.

\begin{enumerate}[leftmargin=*,label=(\roman*)]
\item \emph{Best design, naive}: reject if $\max_k t_k>1.645$. This is what happens when the design is
chosen after looking at the results.
\item \emph{Bootstrap max-$t$ (Reality Check)}: reject if $\max_k t_k$ exceeds the 95\% quantile of the
bootstrap distribution of $\max_k t^*_k$, a studentised version of White's Reality Check
\parencite{white2000}; the step-down extension of \textcite{romano2005} identifies which designs are
significant.
\item \emph{Bootstrap average-$t$}: reject if $\bar t=K^{-1}\sum_k t_k$ exceeds the 95\% quantile of the
bootstrap distribution of $\bar t^*$, following the average-over-windows statistic of
\textcite{rossi2012}.
\end{enumerate}

The bootstrap is the stationary bootstrap of \textcite{politis1994}, with the mean block length chosen
by the method of \textcite{politis2004} and \textcite{patton2009}. The time index is resampled jointly for
all designs, so that each resample contains the same days for every design and the dependence between
designs is preserved; this is how (ii) and (iii) learn from the data how much the $K$ designs have in
common. Bootstrap statistics are recentred, $t^*_k=(\bar d^{*}_k-\bar d_k)/(\hat\sigma_k/\sqrt n)$, so
that they mimic the boundary of the null in which no design has an edge, and are studentised with the
original $\hat\sigma_k$ rather than a bootstrap standard error, a shortcut that makes both tests
somewhat conservative against buy-and-hold when volatility clusters (Section~\ref{sec:warp}).

Both tests address the same intersection null, that no design has an edge; they differ in how the
evidence is weighted. The maximum asks whether at least one design works, the claim implicit in reporting
the best design. The average cannot be driven by one lucky design, although one strong design can carry
it, so rejecting with it does not show that every design works; when the design is a nuisance choice it
is the natural summary and is more stable than any single design. We recommend the average-$t$ test as
the headline and the max-$t$ test as the support for any claim about a particular design.

\paragraph{Why not Bonferroni or the Deflated Sharpe Ratio} Both correct for a stated number of trials:
Bonferroni rejects if $\max_k t_k>z_{1-0.05/K}$ \parencite{dunn1961}, and the Deflated Sharpe Ratio (DSR)
of \textcite{bailey2014dsr} compares the best Sharpe ratio with the maximum expected from $N$
\emph{independent} unskilled trials. With dependent trials the count is replaced by an effective number,
and there is no agreed way to compute one. The genetics literature offers several counts based on the
eigenvalues $\lambda_i$ of the trials' correlation matrix \parencite{nyholt2004,li2005,galwey2009};
\textcite{lopezdeprado2019} clusters the trials; and the participation ratio
$(\sum_i\lambda_i)^2/\sum_i\lambda_i^2$ is a common summary of how many dimensions the trials span. For
our 56 designs these counts range from \NumNeffD{} to about \NumNeffNyholtRound{}
(Section~\ref{sec:dependence}), and a correction is only as good as the count it is fed. We report
Bonferroni and the DSR in Table~\ref{tab:size} for reference, with $N=56$ and with eigenvalue counts
estimated within each simulated market; the bootstrap tests need no count.

% =====================================================================================
\section{Simulation design}\label{sec:design}

\subsection{Rules}

\begin{table}[tbp]
\centering
\caption{The \NumRules{} trading rules: four of the five families of \textcite{sullivan1999}, with the
parameter grid of each family.}
\label{tab:rules}
\input{\tabpath/tab_rules.tex}
\end{table}

The rule universe (Table~\ref{tab:rules}) contains $J=\NumRules$ rules from four of the five families of
\textcite{sullivan1999}: filter rules, moving-average crossovers, support-and-resistance rules and channel
breakouts (the fifth, on-balance volume, needs volume data that index series lack). Each family's grid
is a subset of theirs, spanning the same ranges with a few hundred rules rather than several thousand;
definitions are in Appendix~\ref{app:impl}. Rules use closing prices of the total-return index, since no
price index exists for the long US series; the dividend step this adds to every signal is negligible at
daily frequency. Short filter and breakout rules behave like short-term momentum rules; long moving
averages follow trends that last months. A first version of this study used a smaller universe of
\NumOldRules{} moving-average and channel rules; it was replaced by the present one to bring the rules
close to those of \textcite{sullivan1999}, with the designs and the tests unchanged
(Section~\ref{sec:limits} gives the result under the first universe).

\subsection{Designs}

\begin{table}[tbp]
\centering
\caption{The 56 walk-forward designs: four test-block lengths, two gaps and seven training-window
settings.}
\label{tab:designs}
\input{\tabpath/tab_designs.tex}
\end{table}

Table~\ref{tab:designs} lists the grid. Training lengths of one, two and four years and test blocks of a
month to a year span the ranges used in practice; growing-rolling windows grow by half a test block per
block, and the gap is zero or one month. Testing starts on the same day for all designs, after a
\NumBurnIn-day burn-in (the longest rule look-back) plus the longest training window and gap. The grid
was fixed before any result was computed and has not changed since. Two other parts of the study did
change after first results, the rule universe (above) and the range of the alternatives (below); both
changes are recorded in the decision log of the replication package.

\subsection{Markets with and without predictability}

Each simulated market has $T=\NumT$ trading days, about 21 years, of which the last \NumNoos{} are out of
sample. Two null processes have no predictability: \emph{GARCH}, $r_t=\sigma_t z_t$ with
$\sigma^2_t=\omega+\alpha r^2_{t-1}+\beta\sigma^2_{t-1}$ and $z_t$ standardised Student-$t$
\parencite{engle1982,bollerslev1986,bollerslev1987}, estimated by maximum likelihood on daily US market
returns from July 1990 to September 2026 ($\omega=\NumGarchOmega$, $\alpha=\NumGarchAlpha$,
$\beta=\NumGarchBeta$, $\nu=\NumGarchNu$); and \emph{i.i.d.}\ Student-$t$ returns with the same variance
and tails. Daily shocks are truncated at $\pm\NumCapPct$ so that a simulated price cannot become
negative; the truncation binds on \NumCappedDays{} days in \NumCappedReps{} of the \NumNullReps{} GARCH
markets. A block bootstrap of historical returns is not a null, because it keeps the short-run
autocorrelation that short breakout rules exploit.

Both nulls have a zero mean excess return and zero autocorrelation. Against cash a zero mean is the
boundary of the null. Against buy-and-hold it is the least favourable case: with a positive premium, a
strategy that is sometimes flat or short gives up part of the premium and every procedure rejects less
often. With a premium of \NumPremium{} a year added to the GARCH null (\NumPremiumReps{} markets), a
single design rejects in \NumSizeSinglePremium\% of markets and the naive best of 56 in
\NumSizeBestNaivePremium\%. The simulation therefore answers ``how often does design snooping produce a
false positive in a market with nothing to find and no premium to give up'', which is an upper bound for
tests against buy-and-hold in a market with a premium, and not the question the empirical section faces,
where the index carries autocorrelation that is real but not tradable.

Two alternatives plant predictability on top of the GARCH noise $\varepsilon_t$: short-run momentum,
$r_t=\phi\, r_{t-1}+\varepsilon_t$ with $\phi\in\{\NumPhiList\}$, autocorrelations of a size found in
index returns partly because of stale prices \parencite{lo1990nonsync}; and persistent trends,
$r_t=s_t\,\mu/252+\varepsilon_t$ with $s_t\in\{-1,+1\}$ a Markov regime of average duration 126 days
\parencite{hamilton1989} and annual drift $\mu\in\{\NumMuList\}$. The ranges of $\phi$ and $\mu$ were
widened after a pilot with smaller values showed power too low to separate the procedures; the pilot
concerned power only and is recorded in the decision log.

\subsection{Monte Carlo with a bootstrap inside}\label{sec:warp}

Procedures (ii) and (iii) need a bootstrap in every simulated market. We use the warp-speed method of
\textcite{giacomini2013}: each replication draws one joint bootstrap resample, and the bootstrap critical
value is the 95\% quantile, across replications, of the single bootstrap statistic. As a check, the full
bootstrap (\NumFullbootB{} draws) on \NumFullbootReps{} replications of each null gives rejection rates
against buy-and-hold of \NumFullbootMax\% (max-$t$) and \NumFullbootAvg\% (average-$t$) under the GARCH
null and \NumFullbootMaxIid\% and \NumFullbootAvgIid\% under the i.i.d.\ null; against cash the rates
are \NumFullbootMaxCash\% and \NumFullbootAvgCash\% under the GARCH null and \NumFullbootMaxCashIid\%
and \NumFullbootAvgCashIid\% under the i.i.d.\ null. With \NumFullbootReps{} replications each rate has a
Monte Carlo standard error of about \NumFullbootSE{} points. No full-bootstrap rate exceeds 5\%, and
against buy-and-hold the tests reject less often under the GARCH null than under the i.i.d.\ null:
studentising with the original standard error makes them conservative when volatility clusters. The same
check for power, on \NumFullbootPowerReps{}
replications of each of the two alternatives quoted in Section~\ref{sec:power}, gives
\NumFullbootPowerMax\% (max-$t$) and \NumFullbootPowerAvg\% (average-$t$) for momentum and
\NumFullbootPowerMaxMu\% and \NumFullbootPowerAvgMu\% for trends, close to the warp-speed estimates
reported there. Rates for $K<56$ average over \NumSubsets{} random subsets of $K$ designs. Nulls use
\NumNullReps{} replications and each alternative \NumAltReps.

% =====================================================================================
\section{Simulation results}\label{sec:mc}

\subsection{How many designs are there, really?}\label{sec:dependence}

\begin{figure}[tbp]
\centering
\includegraphics[width=\textwidth]{\figpath/fig_corr.pdf}
\caption{Dependence between the 56 designs under the GARCH null, with buy-and-hold as the benchmark:
correlation across \NumNullReps{} simulated markets of the design $t$-statistics. Designs are ordered as
in Table~\ref{tab:designs}: by test length $h$, then gap $g$, then scheme and training length. With cash
as the benchmark the mean correlation falls from \NumMeanCorrT{} to \NumMeanCorrTCash.}
\label{fig:corr}
\end{figure}

The 56 designs are far from 56 independent experiments, for two reasons. The first is \emph{shared
selection}: they choose among the same rules on overlapping training windows, and designs with similar
windows hold the same rule most of the time (under the null a design selects on average \NumMCdistinct{}
distinct rules over the span). The second is \emph{shared benchmark noise}: the statistic of every
design is computed from $d^{\theta}_t=(p^{\theta}_{t-1}-1)\,e_t$, and the term $-e_t$ is common to all
designs.\footnote{If the positions of two designs were independent of each other and of returns, with
mean zero, and returns had mean zero, the correlation between their series $d_t$ would be $1/(1+\E[p^2])\ge 1/2$: the covariance is
$\E[(p-1)(p'-1)]\,\E[e^2]=\E[e^2]$ and each variance is $(1+\E[p^2])\,\E[e^2]$. Shared selection adds to
this floor.} With cash as the benchmark the common term is absent and only the first source remains.

Figure~\ref{fig:corr} shows the correlation, across simulated markets, of the 56 design $t$-statistics
under the GARCH null against buy-and-hold: the mean pairwise correlation is \NumMeanCorrT{} and the
smallest \NumMinCorrT. Against cash the mean is \NumMeanCorrTCash{} and the smallest \NumMinCorrTCash.
One way to turn the dependence into a number asks how many independent 5\% tests would have to be run
for their maximum to reject as often as the naive best of 56, $\log(1-\mathrm{FPR})/\log 0.95$. With the
false-positive rate of \NumSizeBestNaive\% found below against buy-and-hold the answer is \NumNindep;
with the \NumSizeBestNaiveCash\% found against cash it is \NumNindepCash. For this grid, ``the best of
56 designs'' is the best of about \NumNindepRound{} independent tests against buy-and-hold and about
\NumNindepCashRound{} against cash.

That number is known only from the null simulation. Counts that can be computed from one sample use the
eigenvalues of the correlation matrix of the 56 daily series $d^{\theta}_t$ (Appendix~\ref{app:impl}).
Averaged over the simulated markets, with buy-and-hold as the benchmark, the mean correlation of the
daily series is \NumMeanCorrD{} and the participation ratio is \NumNeffD, the count of
\textcite{galwey2009} \NumNeffGalwey, that of \textcite{li2005} \NumNeffLiji{} and that of
\textcite{nyholt2004} \NumNeffNyholt. The four counts differ by an order of magnitude, and none of them
is designed to match the independent-tests number.

\subsection{False positives}

\begin{figure}[tbp]
\centering
\includegraphics[width=\textwidth]{\figpath/fig_size.pdf}
\caption{False-positive rates when no rule has skill (GARCH null, \NumNullReps{} simulated markets). (a)
Each dot is one walk-forward design used on its own, tested against buy-and-hold (note the narrow axis).
(b) Reporting the best of $K$ designs inflates the rate as $K$ grows, more against cash than against
buy-and-hold, while the bootstrap tests across designs (shown against buy-and-hold) keep it close to
5\%.}
\label{fig:size}
\end{figure}

\begin{table}[tbp]
\centering
\caption{False-positive rates under the two null processes, $K=56$ designs, against buy-and-hold and
against cash.}
\label{tab:size}
\input{\tabpath/tab_size.tex}
\end{table}

A design fixed in advance has the size it claims (Figure~\ref{fig:size}a): against buy-and-hold the
false-positive rate of every design lies between \NumSizeSingleMin\% and \NumSizeSingleMax\%, with an
average of \NumSizeSingle\%, and against cash the average is \NumSizeSingleCash\%. Choosing the best of
\NumRules{} rules inside every training window does not contaminate the out-of-sample test, because the
test only uses days the selection never saw.

The picture changes as soon as the design is chosen after the fact (Figure~\ref{fig:size}b). Against
buy-and-hold the best of two designs has a false-positive rate of \NumSizeBestTwo\%, the best of eight
\NumSizeBestEight\%, and the best of all 56 \NumSizeBestNaive\%. Against cash the best of 56 rejects in
\NumSizeBestNaiveCash\% of null markets (Table~\ref{tab:size}, Panel~B). Both rates rise far more slowly
than for independent tests, where 56 tries would give $1-0.95^{56}\approx94\%$, and the gap between them
is the shared benchmark noise of Section~\ref{sec:dependence}. For comparison, picking the best of
\NumRules{} rules over the whole sample without walk-forward analysis, the classic snooping exercise, has
a false-positive rate of \NumSizeRulesNaive\% against buy-and-hold, which the Reality Check brings back
to \NumSizeRulesRC\%.

The bootstrap max-$t$ and average-$t$ tests across designs reject in \NumSizeMaxT\% and \NumSizeAvgT\%
of null markets against buy-and-hold and in \NumSizeMaxTCash\% and \NumSizeAvgTCash\% against cash, and
the results are similar under the i.i.d.\ null. The one rate above 5\% by more than Monte Carlo error is
that of the average-$t$ test against cash under the GARCH null; it comes from the warp-speed shortcut,
since the full bootstrap of Section~\ref{sec:warp} rejects in \NumFullbootAvgCash\% of markets in the
same case. Corrections that take a count of trials depend on the
count (Table~\ref{tab:size}, Panel~D). Bonferroni fed with $N=56$ rejects in \NumSizeBonf\% of null
markets; with the counts of Nyholt, Li and Ji, and Galwey in \NumSizeBonfNyholt\%, \NumSizeBonfLiji\%
and \NumSizeBonfGalwey\%; and with the participation ratio in \NumSizeBonfEff\%. The DSR rejects in
\NumSizeDSR\% with $N=56$ and in \NumSizeDSREff\% with the participation ratio. The published
eigenvalue counts are safe but conservative, and the participation ratio, which counts the
dimensions the designs span and not the independent chances their maximum gets, is far too liberal. The
count that makes Bonferroni close to exact is the independent-tests number of
Section~\ref{sec:dependence}, about \NumNindep{} (rejection rate \NumSizeBonfIndep\%), which is known
only from the null simulation and changes with the benchmark. The bootstrap tests reach the same place
without a count.

\subsection{Power}\label{sec:power}

\begin{figure}[tbp]
\centering
\includegraphics[width=\textwidth]{\figpath/fig_power.pdf}
\caption{Rejection rates against buy-and-hold when predictability is planted, $K=56$ designs,
\NumNoos{} out-of-sample days per market. (a) Short-run momentum with AR(1) coefficient $\phi$. (b) Persistent trends with annual
drift $\mu$. The grey band is the range of single designs and the dashed line their median. The naive
best-of-56 rule is shown with a size-adjusted critical value (the 95\% quantile of its statistic under
the GARCH null); with the nominal threshold it rejects more often, at a false-positive rate of
\NumSizeBestNaive\%.}
\label{fig:power}
\end{figure}

Figure~\ref{fig:power} reports how often each procedure detects predictability that is there, with
buy-and-hold as the benchmark. The comparison a reader will make is with the single design they would have used, and the answer depends on
which one. Against the median design the robust tests lose nothing and in fact gain: with momentum of
$\phi=\NumPowerPhi$
the average-$t$ test rejects in \NumPowerAvg\% of markets and the Reality Check in \NumPowerMax\%,
against \NumPowerMedian\% for the median design; with trends of $\mu=\NumPowerMu$ a year the rates are
\NumPowerAvgMu\%, \NumPowerMaxMu\% and \NumPowerMedianMu\%. Against the best design they do cost power:
at $\phi=\NumPowerPhi$ single designs range from \NumPowerSingleMin\% to \NumPowerSingleMax\%, the
expanding-window designs, which a researcher could reasonably have fixed in advance, average
\NumPowerExpanding\%, and the robust tests give up about \NumPowerGapPoints{} percentage points against
them. That dispersion is the point: committing to one design is a gamble on having picked a good one,
and picking it after the fact costs a false-positive rate of \NumSizeBestNaive\%. The size-adjusted
naive rule, with the oracle critical value from the null simulation, rejects in \NumPowerNaiveAdj\% of
markets at $\phi=\NumPowerPhi$, essentially the Reality Check, whose bootstrap estimates that critical
value from the data at hand. Bonferroni and the DSR with the participation-ratio count reject more often
still, but at false-positive rates of \NumSizeBonfEff\% and \NumSizeDSREff\% those rates are not power.

\subsection{The in-sample to out-of-sample gap}

Practitioners often compare the return of the selected rules in the training windows with their return
in the test blocks and read a small gap as evidence of skill. Under the GARCH null the selected rules
earn \NumISann\% a year in sample on average and \NumOOSann\% out of sample: a large in-sample return
is what selection produces in a market with nothing to find, and a drop to zero out of sample is the
expected outcome. Only a test against a benchmark under the appropriate null tells skill from luck
\parencite{inoue2005,harvey2015}.

% =====================================================================================
\section{Empirical illustration}\label{sec:empirical}

\subsection{Data}

\begin{table}[tbp]
\centering
\caption{Data: the seven return series, their samples and the buy-and-hold record over the
out-of-sample span.}
\label{tab:data}
\input{\tabpath/tab_data.tex}
\end{table}

The main series is the daily total return of the US value-weighted stock market from July 1926 to
September 2026, with the one-month Treasury-bill rate, from Kenneth French's data library
\parencite{french2026}, the closest public counterpart of the S\&P~500 over a century. We add five
regional series from the same library from July 1990 and, as a short recent check, the S\&P~500 price
index from FRED \parencite{fred2026}; Table~\ref{tab:data} summarises them. Costs are \NumCostBp{} basis
points per unit of position change (0.2\% for a switch from long to short), a level that is realistic
for index futures today and flatters the early decades, when trading the underlying stocks cost several
times as much \parencite{jones2002}. The series illustrates the tests; it is not a test of technical
analysis, because a value-weighted index of all listed stocks could not be traded before index futures
began in 1982, and the strategies change position about \NumUSturnover{} times a year.

\subsection{Results}

\begin{table}[tbp]
\centering
\caption{Walk-forward strategies against buy-and-hold, 56 designs per series: after costs, before costs,
and before costs with every trade executed one day after its signal.}
\label{tab:empirical}
\input{\tabpath/tab_empirical.tex}
\end{table}

\begin{figure}[tbp]
\centering
\includegraphics[width=\textwidth]{\figpath/fig_us_fan.pdf}
\caption{Growth of \$1 invested above T-bills in the US market, \NumUSoosFirst--2026: buy-and-hold
(black) and the 56 walk-forward strategies (one line per design), (a) after a cost of 10 basis points per
unit of position change and (b) before costs. The shaded area marks the years before index futures
existed.}
\label{fig:fan}
\end{figure}

After costs the answer is nearly uniform (Table~\ref{tab:empirical}, Panel~A; Figure~\ref{fig:fan}a). In
the US series no design beats buy-and-hold: the best $t$-statistic is \NumUSnetBestT{} and the mean is
\NumUSnetMeanT, and the design-robust tests are nowhere near rejecting ($p$ = \NumUSnetPmax{} for the
max-$t$ test, $p$ = \NumUSnetPavg{} for the average-$t$). The same holds in every regional series except
Asia Pacific, where \NumAPnetSig{} of 56 designs are significant after costs and both tests reject ($p$ =
\NumAPnetPmax{} and \NumAPnetPavg). That index has a first-order autocorrelation of \NumAPac, which comes
from markets that close at different hours and is a property of the index, not of any portfolio an
investor could trade at closing prices.

Before costs the US evidence looks strong (Panel~B): \NumUSgrossSig{} of 56 designs beat buy-and-hold at
the 5\% level, the best has $t=\NumUSgrossBestT$, and the design-robust tests reject ($p$ =
\NumUSgrossPmax{} and \NumUSgrossPavg). This is the situation the tests are for. The rejection says that
the result is not an artefact of choosing the design; it does not say what the 56 designs have found.

They have not found 56 things. The mean correlation of the 56 daily series $d^{\theta}_t$ is
\NumUSmeanCorr{} and their participation ratio \NumUSneff; \NumUSdistinctLow{} of the 56 designs select
fewer than ten distinct rules over the whole span, and one rule, \NumUStopRule, wins \NumUStopShare\% of
all training windows across all designs (the three most selected rules win \NumUStopThreeShare\%). The
full-sample best of the \NumRules{} rules, without walk-forward selection, is \NumUSclassicRule{} with
$t=\NumUSclassicT$ and a Reality Check $p$-value of \NumUSclassicPrc, so the gross result is not specific
to walk-forward selection either.

Nor could the profits have been earned. A rule that acts on a move of 0.5\% trades the first-order
autocorrelation of the index, which was \NumACpre{} before 1982 against \NumACpost{} since.
Autocorrelation of that size in a broad value-weighted index is what nonsynchronous trading produces:
stale prices of stocks that did not trade enter the index, which absorbs news with a delay that looks
like predictability but cannot be captured by trading at the quoted close \parencite{lo1990nonsync}. The
direct check is to execute every trade one day after its signal \parencite{bessembinder1998}, in the
training windows and in the test blocks alike (Panel~C). With that single change no design is
significant in any of the seven series. In the US series the number of significant designs
falls from \NumUSgrossSig{} to \NumUSlagSig{} and the best $t$-statistic from \NumUSgrossBestT{} to
\NumUSlagBestT; in Asia Pacific the count falls from \NumAPgrossSig{} to \NumAPlagSig{} and in the
developed-markets index from \NumDVgrossSig{} to \NumDVlagSig. \textcite{day2002} and
\textcite{ready2002} showed that the same mechanism, with dividends, accounts for much of the profit of
the \textcite{brock1992} rules. The gross profits are also concentrated in the decades before index
futures existed (Figure~\ref{fig:fan}b; Appendix~\ref{app:eras} gives the split by period), which matches
the rise and fall of technical trading profits documented for US indices
\parencite{sullivan1999,taylor2014,zakamulin2014}.

The last column of Table~\ref{tab:empirical}, the probability of backtest overfitting (PBO) of
\textcite{bailey2017}, asks how often the design that is best in one half of the sample ranks below the
median in the other half: \NumUSnetPBO\% of splits in the US series after costs and \NumAPnetPBO\% in
the Asia Pacific series. The low values say that the ranking of designs is stable across halves, as
designs this dependent must be; PBO is a statement about the choice among designs, not about whether any
of them beats the benchmark.

% =====================================================================================
\section{Discussion}\label{sec:discussion}

\subsection{Recommendations}

\begin{enumerate}[leftmargin=*]
\item \textbf{Check the accounting.} Use a splitter in which test blocks tile the evaluation span and all
designs share it; report the number of out-of-sample days.
\item \textbf{Fix the design in advance, or report all of them}, with the benchmark and the correlation
between the designs' statistics. The cost of choosing among designs depends on both.
\item \textbf{Test across designs.} Report the bootstrap average-$t$ test as the headline and the Reality
Check (max-$t$) for the claim that at least one design works; both take seconds on the per-design results.
\item \textbf{Do not lean on a trial count.} Corrections that take one, such as Bonferroni or the DSR,
are too conservative with the nominal number of designs, and effective numbers computed from the same
correlation matrix differ by an order of magnitude. Where one is used, the designs count as trials
alongside the rules and parameters, and the count and its formula should be reported with the result.
\item \textbf{Compare with a benchmark, after costs, on a tradable series}, with dependence-robust
standard errors, and repeat the test with every trade executed one day after its signal. A result that
disappears with the lag is index autocorrelation, not a trading profit.
\item \textbf{Do not read the in-sample to out-of-sample gap as evidence.} Selection produces a large
in-sample return even when there is nothing to find.
\end{enumerate}

\subsection{Limitations and further work}\label{sec:limits}

The simulations use one selection criterion and one rule universe, and the level of the inflation
depends on the universe: with the \NumOldRules{} rules of the first version of this study the naive best
of 56 designs rejected in \NumOldSizeBestNaive\% of GARCH markets against buy-and-hold, against
\NumSizeBestNaive\% with the present \NumRules. The mechanism and the ranking of the procedures were the
same. The null processes have no autocorrelation and no structural breaks; under breaks the choice
between rolling and expanding windows also affects power \parencite{pesaran2007,inoue2017}. The
simulated markets have \NumNoos{} out-of-sample days, about \NumNoosShareUS\% of the US sample, so the
power figures are for a \NumNoosYears-year record, and power is reported against buy-and-hold only.

The size of the two tests is established here by simulation under two null processes, not by a proof.
The theory of the Reality Check and its step-down version \parencite{white2000,romano2005} covers
stationary return differences, and a strategy whose training window grows is not exactly stationary.
Against buy-and-hold both tests are conservative under volatility clustering, because the bootstrap
statistics are studentised with the original standard error and one block length serves all designs. The tests compare
mean returns, not Sharpe ratios, although a strategy that is often flat has a lower volatility than
buy-and-hold. Rules are selected on returns before costs; selecting on returns after costs would favour
slower rules.

We test across walk-forward designs rather than switching to combinatorial purged cross-validation
\parencite{lopezdeprado2018,arian2024}, because walk-forward paths are tradable histories and are what
practitioners report; the two approaches answer different questions and can be used together. The
natural next step is a wider grid in which the selection criterion, the rule universe, the benchmark, the
cost charged inside selection and the start of the evaluation span vary jointly with the window settings;
the formulation and the tests apply unchanged, and the dependence between the designs of such a grid
would show how fast the inflation grows with what is varied.

% =====================================================================================
\section{Conclusions}\label{sec:conclusion}

Walk-forward analysis fixes rule-level snooping but not design-level snooping. Within a family in which
every design is scored on the same days, a design fixed in advance has the false-positive rate it claims,
and reporting the best of 56 designs raises it from 5\% to about \NumSizeBestNaive\% when the benchmark
is buy-and-hold and to about \NumSizeBestNaiveCash\% when it is cash. The designs are dependent because
they select among the same rules and, against buy-and-hold, because they share the benchmark's noise;
the best of 56 is the best of about \NumNindepRound{} independent tests in the first case and about
\NumNindepCashRound{} in the second. These numbers are properties of this grid, not of walk-forward
analysis: a grid that also varied the selection criterion, the rule universe, the cost model or the start
of the evaluation span would contain designs that disagree more, and the inflation would grow
accordingly.

The remedy does not depend on the grid or on a count of trials. Two bootstrap tests that treat the
designs the way the Reality Check treats rules, as a family of hypotheses tested jointly with a bootstrap
that preserves their dependence, keep the false-positive rate close to 5\% under both benchmarks and,
against buy-and-hold, keep the power of a typical design. A Bonferroni correction gives rejection rates
between \NumSizeBonf\% and \NumSizeBonfEff\% on the same designs, depending on the count of trials it is
fed.

On a century of US data the tests agree with every single design that nothing survives costs. Before
costs they reject, and the right response is to ask what the designs have in common: here a
short-horizon rule that trades index autocorrelation, whose profit disappears when trades are executed
one day late. The tests remove the design as a source of false discoveries; they do not remove the need
to check what the data can be traded at. The same accounting and the same tests can be applied to other
models that are re-estimated over time and evaluated with a walk-forward design.

\section*{Declarations}
\paragraph{Data and code availability} All data are public \parencite{french2026,fred2026}. A replication
package that reproduces every number, table and figure with a single command, with fixed random seeds, is
available at \url{https://github.com/ccastro6/how_many_ways_to_walkforward} and archived at
\url{https://doi.org/10.5281/zenodo.24237608}.

\paragraph{Funding and conflicts of interest} This work received no specific funding. The authors declare
no conflicts of interest.

\paragraph{Use of generative AI} The authors used Claude Opus 5.5 (Anthropic) to assist with literature
searching, coding and language editing, and take full responsibility for the content.


{\emergencystretch=1.5em
\printbibliography[heading=bibintoc,title={References}]}

\appendix
\section{Implementation details}\label{app:impl}

\paragraph{Rules} Filter rules work on log prices: a rule with filter $x$ goes long when the price has
risen $x$ above its reference low and short when it has fallen $x$ below its reference high, reversing at
each signal; the reference extremes run from the last signal. Variants liquidate a position after an
adverse move $y<x$ (the next entry then needs a full move of $x$), hold each position for $c$ days and
then go flat, or take the reference extremes over the previous $e$ days. Moving-average rules are long
when the $n$-day average exceeds the $m$-day average and short otherwise; with a band $b$ they are flat
while the two averages are within $b$ of each other. Support-and-resistance rules go long when the close
exceeds the maximum of the previous $n$ closes and short when it falls below the minimum, by more than a
band $b$ where one applies; the position is held until the opposite signal or, in the holding-period
variants, for $c$ days and then closed. Channel breakouts apply the same breakout signals only when a
channel exists, that is, when the maximum of the previous $n$ closes is within $x$ of the minimum, and
hold the position $c$ days. The parameter grids are those of Table~\ref{tab:rules}.

\paragraph{Sharpe ratios and ties} In-sample Sharpe ratios are computed for all rules and all windows in
constant time per window from cumulative sums of returns and squared returns, taken on returns centred
at each rule's full-sample mean so that the variance loses no precision. Rules with identical returns
over a window tie exactly; to make the arg-max independent of floating-point noise in the cumulative sums
the ratios are rounded to ten decimals before it, and ties are broken by rule order. Ties are common for
short windows, because several rules can take identical positions over a year: in the US sample
\NumUStiedWindows{} of the \NumUSwindows{} training windows across the 56 designs (\NumUStiedShare\%)
have a tied maximum, up to \NumUSmaxTied{} of \NumUSmaxTiedBlocks{} windows in design
\NumUSmaxTiedDesign{} (\NumUSmaxTiedDesc). Under the GARCH null
\NumMCtiedShare\% of training windows are tied.

\paragraph{Turnover and costs} Position changes are counted on every day, including the day a design
switches from one rule to another at a block boundary. The first position of the evaluation span counts
as a change from flat.

\paragraph{Bootstrap} The stationary bootstrap uses circular wrapping and the Politis--White mean block
length computed on the cross-design average of $d_t$. Empirical $p$-values use 1,999 draws and the
convention $(1+\#\{\text{exceedances}\})/(1+B)$. With cash as the benchmark the same steps are applied to
the strategy returns $y^{\theta}_t$, with a block length computed on their cross-design average and a
separate bootstrap draw.

\paragraph{Execution lag} In the lagged runs of Section~\ref{sec:empirical} the position decided at the
close of day $t$ is entered at the close of day $t+1$, so that rule $j$ earns $p_{j,t-2}\,e_t$ on day
$t$. The lag applies to the returns used for selection in the training windows as well as to the returns
earned in the test blocks; everything else is unchanged.

\paragraph{Effective number of designs} Let $\lambda_1,\dots,\lambda_K$ be the eigenvalues of the
correlation matrix of the design statistics (across simulated markets) or of the daily series
$d^{\theta}_t$ (within one sample). The participation ratio is $(\sum_i\lambda_i)^2/\sum_i\lambda_i^2$,
which equals $K$ for independent designs and 1 for identical ones. The count of \textcite{nyholt2004} is
$1+(K-1)\,(1-\mathrm{Var}(\lambda)/K)$, with $\mathrm{Var}(\lambda)$ the sample variance of the
eigenvalues; that of \textcite{li2005} is $\sum_i\bigl(\mathbf{1}\{\lambda_i\ge 1\}+\lambda_i-\lfloor\lambda_i\rfloor\bigr)$;
and that of \textcite{galwey2009} is $(\sum_i\sqrt{\lambda_i})^2/\sum_i\lambda_i$, with negative eigenvalues
set to zero. The independent-tests count is $\log(1-\mathrm{FPR})/\log(1-0.05)$, where FPR is the
false-positive rate of the naive best-of-$K$ rule under the null. Bonferroni with an effective number
$N$ rejects when $\max_k t_k>z_{1-0.05/N}$. The DSR with an effective number uses the expression of
\textcite{bailey2014dsr} with per-period Sharpe ratios, the cross-design variance of Sharpe ratios, the
skewness and kurtosis of the best design's returns over buy-and-hold, and $N$ equal to the within-sample
participation ratio.

\paragraph{Probability of backtest overfitting} PBO uses combinatorially symmetric cross-validation with
16 blocks (12,870 splits) applied to the matrix of daily differences $d^{\theta}_t$ across the 56 designs
\parencite{bailey2017}.

\section{US results by period}\label{app:eras}

\begin{figure}[tbp]
\centering
\includegraphics[width=\textwidth]{\figpath/fig_eras.pdf}
\caption{$t$-statistics of the 56 walk-forward designs against buy-and-hold in three periods of the US
sample, before and after costs. Each dot is a design. Designs that converge on the same rule produce
identical statistics and appear as horizontal rows.}
\label{fig:eras}
\end{figure}

\begin{table}[tbp]
\centering
\caption{US market by period: the 56 designs before costs, after costs, and before costs with every
trade executed one day after its signal.}
\label{tab:eras}
\input{\tabpath/tab_eras.tex}
\end{table}

The strategies are run once over the full sample and their out-of-sample days are split into
\NumUSoosFirst--1959, 1960--1989 and 1990--2026; the split follows the break around 1990 reported in the
literature and is a description of the record, not a test designed in advance
(Figure~\ref{fig:eras}, Table~\ref{tab:eras}). Before 1960, \NumEraOneSig{} of 56 designs are
significant before costs and \NumEraOneNetSig{} after costs; the naive best design rejects while the
design-robust tests do not ($p$ = \NumEraOnePmax{} for the max-$t$ test and \NumEraOnePavg{} for the
average-$t$), the situation in which choosing the design after the fact decides the conclusion. In
1960--1989, \NumEraTwoSig{} of 56 designs are significant before costs and \NumEraTwoNetSig{} after costs
(\NumEraTwoNetStepm{} under the step-down procedure; average-$t$ $p$ = \NumEraTwoNetPavg); the
first-order autocorrelation of the index is \NumACtwo{} in this period, against \NumACone{} before 1960
and \NumACthree{} since 1990. Two checks show that the 1960--1989 result is not a trading profit. With
every trade executed one day after its signal, \NumEraTwoLagSig{} of 56 designs are significant before
costs and the largest $t$-statistic falls from \NumEraTwoMaxT{} to \NumEraTwoLagMaxT. And the average gain
per unit traded in that period, \NumEraTwoBreakeven{} basis points, is the cost at which the average
design breaks even; it is well below what trading the underlying stocks cost at the time
\parencite{jones2002}. After the assumed cost of \NumCostBp{} basis points the gain was complete by 1982:
the strategies' cumulative log return relative to buy-and-hold averages \NumUSnetRelPreFutures{} at the
end of 1981, \NumUSnetRelEndEighties{} at the end of 1989 and \NumUSnetRelEnd{} at the end of the sample.
Since 1990, \NumEraThreeSig{} designs are significant and the largest $t$-statistic is
\NumEraThreeMaxT. The pattern is consistent with predictability being competed away once it is known and
cheap to trade \parencite{lo2004,timmermann2004,mclean2016}, and with the nonsynchronous-price
explanation of the earlier profits \parencite{day2002,ready2002}.

\end{document}
